% ============================================================
% Chapter 13: Problem 12 — Kronecker's Theorem on Abelian Extensions
% ============================================================

\chapter{Kronecker's Theorem on Abelian Extensions}

\begin{solvedbox}
Kronecker's Jugendtraum is \textbf{solved}: class field theory provides a complete description of abelian extensions of number fields. For $\mathbb{Q}$ the Kronecker--Weber theorem says every abelian extension lies inside a cyclotomic field. For imaginary quadratic fields, complex multiplication gives the required special values (elliptic functions at torsion points). For general number fields, the problem is resolved by class field theory, though Hilbert's original vision---generating all abelian extensions by special values of analytic functions---is realized fully only in the two cases above.
\end{solvedbox}

% ------------------------------------------------------------
\section{Historical Context}
% ------------------------------------------------------------

In the 1850s, Leopold Kronecker (1823--1891) became fascinated by a discovery: the abelian extensions of $\mathbb{Q}$---those Galois extensions whose Galois group is abelian---could all be obtained by adjoining roots of unity. The field $\mathbb{Q}(\zeta_n)$, where $\zeta_n = e^{2\pi i / n}$ is a primitive $n$th root of unity, has Galois group $(\mathbb{Z}/n\mathbb{Z})^\times$, which is abelian. Kronecker suspected that a similar phenomenon held for imaginary quadratic fields: their abelian extensions should come from special values of elliptic functions.

He called this vision his \textbf{\emph{Jugendtraum}}---his ``youthful dream.'' The dream was that \emph{every abelian extension of a number field could be generated by special values of certain analytic functions}. For $\mathbb{Q}$, the functions were the exponential function $e^{2\pi i z}$ (giving roots of unity). For imaginary quadratic fields, the functions were elliptic modular functions and values of the $j$-invariant. For general number fields, the analytic functions would be the \emph{automorphic functions} of the appropriate type.

This problem, as Hilbert formulated it, was about finding the full scope of this phenomenon.

\subsection{What Is an Abelian Extension?}

Before we go further, we need a precise definition.

\begin{definition}
A finite extension $L$ of a number field $K$ is called \textbf{Galois} if it is both normal (every irreducible polynomial with a root in $L$ splits completely in $L$) and separable. The \textbf{Galois group} $\Gal(L/K)$ is the group of field automorphisms of $L$ that fix $K$ pointwise. If $\Gal(L/K)$ is an abelian group, then $L/K$ is called an \textbf{abelian extension}.
\end{definition}

The abelian condition is strong. It means that the symmetries of the extension commute with each other---a constraint that turns out to be sufficient to force the structure of the extension to be very special.

\subsection{Cyclotomic Fields as Examples}

The simplest and most important examples of abelian extensions are the \textbf{cyclotomic fields}.

\begin{definition}
Let $n \geq 1$ and let $\zeta_n = e^{2\pi i/n}$. The $n$th \textbf{cyclotomic field} is $\mathbb{Q}(\zeta_n)$.
\end{definition}

Why is the Galois group abelian? Consider an automorphism $\sigma \in \Gal(\mathbb{Q}(\zeta_n)/\mathbb{Q})$. It is determined by where it sends $\zeta_n$, and $\sigma(\zeta_n)$ must also be a primitive $n$th root of unity, so $\sigma(\zeta_n) = \zeta_n^k$ for some $k$ coprime to $n$. The map
\[
\Gal(\mathbb{Q}(\zeta_n)/\mathbb{Q}) \longrightarrow (\mathbb{Z}/n\mathbb{Z})^\times,\qquad
\sigma \mapsto k \pmod{n}
\]
is an isomorphism. The group $(\mathbb{Z}/n\mathbb{Z})^\times$ is abelian, so every cyclotomic extension is abelian.

\section{The Problem}

\begin{problem_statement}[Kronecker's Theorem on Abelian Extensions]
Prove that every finite abelian extension of a number field can be generated by special values of certain analytic functions. In particular:
\begin{enumerate}
    \item For $K = \mathbb{Q}$, every abelian extension is contained in a cyclotomic field $\mathbb{Q}(\zeta_n)$.
    \item For $K$ an imaginary quadratic field, every abelian extension is generated by values of elliptic modular functions.
    \item For general number fields, find the appropriate analytic functions that generate their abelian extensions.
\end{enumerate}
\end{problem_statement}

Hilbert called this problem ``Kronecker's theorem'' because Kronecker had already stated the result for $\mathbb{Q}$ and for imaginary quadratic fields, though without complete proofs. The problem is to prove these claims and extend them to all number fields.

% ------------------------------------------------------------
\section{A Detailed Worked Example: Quadratic Fields and Kronecker--Weber}
% ------------------------------------------------------------

The Kronecker--Weber theorem says that every abelian extension of $\mathbb{Q}$ lies inside a cyclotomic field. Let us make this concrete for the simplest non-trivial case: quadratic extensions.

Every quadratic extension of $\mathbb{Q}$ is of the form $\mathbb{Q}(\sqrt{d})$ where $d$ is a squarefree integer (positive or negative). The Galois group has order $2$, generated by the automorphism $\sqrt{d} \mapsto -\sqrt{d}$. Since groups of order $2$ are abelian, every quadratic extension is abelian. The Kronecker--Weber theorem therefore predicts that each $\mathbb{Q}(\sqrt{d})$ sits inside some cyclotomic field.

\begin{example}
Consider $\mathbb{Q}(\sqrt{2})$. A primitive 8th root of unity is $\zeta_8 = e^{2\pi i/8} = e^{\pi i/4} = \frac{\sqrt{2}}{2} + i\frac{\sqrt{2}}{2}$. Observe that
\[
\zeta_8 + \zeta_8^{-1} = 2\cos(\pi/4) = \sqrt{2},
\]
so $\sqrt{2} \in \mathbb{Q}(\zeta_8)$. Hence $\mathbb{Q}(\sqrt{2}) \subseteq \mathbb{Q}(\zeta_8)$.
\end{example}

\begin{example}
For $\mathbb{Q}(\sqrt{-1})$, we have $\mathbb{Q}(\sqrt{-1}) = \mathbb{Q}(i) = \mathbb{Q}(\zeta_4)$ because $\zeta_4 = i$ is already a primitive 4th root of unity. Thus the quadratic extension $\mathbb{Q}(i)$ is itself a cyclotomic field.
\end{example}

\begin{example}
For $\mathbb{Q}(\sqrt{5})$, note that $\zeta_5 + \zeta_5^{-1} = 2\cos(2\pi/5) = \frac{-1 + \sqrt{5}}{2}$. This gives $\sqrt{5} = 2\zeta_5 + 2\zeta_5^{-1} + 1$, so $\sqrt{5} \in \mathbb{Q}(\zeta_5)$. Indeed $\mathbb{Q}(\sqrt{5}) \subseteq \mathbb{Q}(\zeta_5)$.
\end{example}

\begin{example}
For $\mathbb{Q}(\sqrt{-3})$, we have $\zeta_3 = e^{2\pi i/3} = -\frac12 + i\frac{\sqrt{3}}{2}$, so $\sqrt{-3} = 2\zeta_3 + 1$. Hence $\mathbb{Q}(\sqrt{-3}) \subseteq \mathbb{Q}(\zeta_3)$.
\end{example}

\begin{example}
For $\mathbb{Q}(\sqrt{3})$, we need a larger cyclotomic field. Using $\zeta_{12}$, one can check that
\[
\zeta_{12} + \zeta_{12}^{-1} = 2\cos(\pi/6) = \sqrt{3},
\]
so $\mathbb{Q}(\sqrt{3}) \subseteq \mathbb{Q}(\zeta_{12})$.
\end{example}

These examples illustrate a general pattern. If $d$ is squarefree and odd, then $\mathbb{Q}(\sqrt{d}) \subseteq \mathbb{Q}(\zeta_{|d|})$ when $d \equiv 1 \pmod{4}$, while $\mathbb{Q}(\sqrt{d}) \subseteq \mathbb{Q}(\zeta_{4|d|})$ otherwise. When $d$ is even, a power of $2$ is needed: $\mathbb{Q}(\sqrt{2}) \subseteq \mathbb{Q}(\zeta_8)$, for instance. The smallest $n$ such that $\mathbb{Q}(\sqrt{d}) \subseteq \mathbb{Q}(\zeta_n)$ is called the \textbf{conductor} of the extension.

\section{Key Developments}

The resolution of this problem spans nearly a century and involves some of the significant ideas in number theory.

\subsection{The Kronecker--Weber Theorem}

The first major step was the proof of the case $K = \mathbb{Q}$.

\begin{theorem}[Kronecker--Weber Theorem]
Every finite abelian extension $L/\mathbb{Q}$ is contained in some cyclotomic field $\mathbb{Q}(\zeta_n)$. Equivalently, the maximal abelian extension $\mathbb{Q}^{\text{ab}}$ of $\mathbb{Q}$ is $\bigcup_{n \geq 1} \mathbb{Q}(\zeta_n)$.
\end{theorem}

Kronecker had stated this theorem in 1853, but his proof contained gaps. A complete proof was given by Heinrich Martin Weber in 1886, and later simplified by David Hilbert in 1896. The theorem tells us that the abelian extensions of $\mathbb{Q}$ are completely understood: they are just subfields of cyclotomic fields.

\subsection{Complex Multiplication}

The second part of Kronecker's dream---the case of imaginary quadratic fields---took longer to fully prove. The theory of \textbf{complex multiplication} provides the analytic functions needed.

For an imaginary quadratic field $K = \mathbb{Q}(\sqrt{-d})$ with $d > 0$, one considers the lattice $\Lambda = \mathbb{Z} + \mathbb{Z}\tau$ in the complex plane, where $\tau$ lies in the upper half-plane and generates the ring of integers of $K$. The \textbf{Weierstrass $\wp$-function}
\[
\wp(z; \Lambda) = \frac{1}{z^2} + \sum_{\omega \in \Lambda \setminus \{0\}} \left(\frac{1}{(z - \omega)^2} - \frac{1}{\omega^2}\right)
\]
is a doubly periodic meromorphic function (an \textbf{elliptic function}) with periods $\Lambda$. Together with its derivative, it satisfies the differential equation $(\wp')^2 = 4\wp^3 - g_2\wp - g_3$, which parametrizes an elliptic curve $E: y^2 = 4x^3 - g_2 x - g_3$.

The \textbf{$j$-invariant} $j(\tau)$ is a modular function given by
\[
j(\tau) = 1728\,\frac{g_2(\tau)^3}{g_2(\tau)^3 - 27\,g_3(\tau)^2}
\]
that classifies elliptic curves up to isomorphism. When $\tau$ lies in an imaginary quadratic field $K$, the elliptic curve $E$ is said to have \textbf{complex multiplication} by $K$: its endomorphism ring is larger than $\mathbb{Z}$, containing an order isomorphic to the ring of integers of $K$.

Kronecker's theorem for imaginary quadratic fields states:

\begin{theorem}[Kronecker's Jugendtraum for Imaginary Quadratic Fields]
Let $K$ be an imaginary quadratic field. Then every abelian extension of $K$ is contained in a field generated by:
\begin{enumerate}
    \item Special values of the $j$-function $j(\tau)$, where $\tau \in K$ lies in the upper half-plane, and
    \item Values of the elliptic functions $\wp$ at torsion points of the associated elliptic curve.
\end{enumerate}
\end{theorem}

Concretely, adjoining $j(\tau)$ to $K$ yields the \textbf{Hilbert class field} of $K$---the maximal unramified abelian extension. Adjoining the values of $\wp$ at $N$-torsion points then builds the ray class fields, which cover all abelian extensions of $K$. For example, if $K = \mathbb{Q}(i)$, then $j(i) = 1728$, and the Hilbert class field of $\mathbb{Q}(i)$ is $\mathbb{Q}(i)$ itself (since $\mathbb{Q}(i)$ has class number $1$). If $K = \mathbb{Q}(\sqrt{-5})$, then $j(\sqrt{-5})$ is a non-trivial algebraic integer generating the Hilbert class field.

This was completed primarily by Kronecker himself and later fully proved by Fueter, Hasse, and others in the early twentieth century. The result is notable: the abelian extensions of an imaginary quadratic field arise from the arithmetic of elliptic curves with complex multiplication.

\subsection{Kronecker's Jugendtraum Explained}

Kronecker's Jugendtraum (youthful dream) was ambitious. He envisioned that \emph{every} abelian extension of \emph{every} number field could be generated by special values of analytic functions analogous to how $\mathbb{Q}(\zeta_n)$ is generated by $e^{2\pi i/n}$. In modern terms, he wanted an explicit class field theory---a constructive recipe for building abelian extensions using transcendental methods.

The difficulty lies in finding the right analytic functions. For $\mathbb{Q}$, the exponential function $e^{2\pi i z}$ works because the maximal abelian extension of $\mathbb{Q}$ is the field generated by all roots of unity, which are precisely the values of $e^{2\pi i z}$ at rational arguments. For imaginary quadratic fields, elliptic functions work because the endomorphism ring of the associated elliptic curve is an order in an imaginary quadratic field, providing the algebraic structure needed.

Why is the general case so hard? For a general number field $K$, one needs a theory of automorphic forms on $\GL_n$ (where $n = [K:\mathbb{Q}]$) that play the role of the exponential function ($n=1$) and elliptic functions ($n=2$). The values of these automorphic forms at special points should generate the abelian extensions. This is the essence of Hilbert's formulation: find the analytic functions for each number field.

Class field theory solves the problem in an algebraic sense---it classifies abelian extensions via the Artin map---but does not give the explicit analytic generators Kronecker wanted. Building explicit generators for general number fields remains an active area of research, often called the \textbf{explicit class field theory} problem. For CM fields (totally imaginary quadratic extensions of totally real fields), the theory of \textbf{Shimura varieties} provides a partial generalization.

\subsection{Class Field Theory}

The full solution to Hilbert's twelfth problem came not through finding special analytic functions for each number field, but through the development of \textbf{class field theory}.

\paragraph{Takagi's Existence Theorem.} In the 1920s, Teiji Takagi proved a sweeping result that is the cornerstone of class field theory:

\begin{theorem}[Takagi, 1920]
For every number field $K$, there is a bijective correspondence between its finite abelian extensions $L/K$ and certain subgroups of the ideal group $I_K$ (the \emph{ideal groups} or \emph{ray class groups}). Moreover, this correspondence respects the splitting of primes.
\end{theorem}

\paragraph{The Artin Map.} Emil Artin, in 1927, gave the final form of this correspondence through the \textbf{Artin reciprocity law}.

\begin{theorem}[Artin Reciprocity]
Let $L/K$ be a finite abelian extension. There exists a canonical homomorphism
\[
\left(\frac{L/K}{\cdot}\right) : I_K^{\mathfrak{m}} \longrightarrow \Gal(L/K),
\]
where $I_K^{\mathfrak{m}}$ denotes the group of fractional ideals of $K$ coprime to the conductor $\mathfrak{m}$. This map is surjective, and its kernel consists precisely of the principal ideals whose generators satisfy certain congruence conditions modulo $\mathfrak{m}$.
\end{theorem}

\subsection{The Artin Reciprocity Map in Concrete Terms}

The Artin map is the bridge between arithmetic (ideals) and symmetry (Galois groups). Let us see what it does in a concrete setting.

For a quadratic extension $\mathbb{Q}(\sqrt{d})/\mathbb{Q}$, the Galois group has two elements: the identity and the map $\sqrt{d} \mapsto -\sqrt{d}$. The Artin map assigns to each prime $p$ (not dividing the discriminant) a Frobenius element $\Frob_p \in \Gal(\mathbb{Q}(\sqrt{d})/\mathbb{Q})$. This Frobenius element encodes how the prime $p$ behaves in the extension:
\[
\Frob_p = \begin{cases}
\text{identity} & \text{if $p$ splits in $\mathbb{Q}(\sqrt{d})$},\\
\text{non-trivial} & \text{if $p$ remains inert}.
\end{cases}
\]

Explicitly, the Artin map sends $p$ to the automorphism $\sqrt{d} \mapsto \left(\frac{d}{p}\right)\sqrt{d}$, where $\left(\frac{d}{p}\right)$ is the \textbf{Legendre symbol}. In other words,
\[
\left(\frac{\mathbb{Q}(\sqrt{d})/\mathbb{Q}}{p}\right) : \sqrt{d} \longmapsto \left(\frac{d}{p}\right)\sqrt{d}.
\]

This is a beautiful concrete instance of Artin reciprocity: the splitting behavior of a prime in a quadratic extension is completely determined by a congruence condition modulo $d$, captured by the Legendre symbol. Quadratic reciprocity is then equivalent to the statement that the Artin map factors through the ray class group of conductor $d$ (or $4d$).

For cyclotomic fields $\mathbb{Q}(\zeta_n)/\mathbb{Q}$, the Artin map is even simpler: a prime $p \nmid n$ maps to the automorphism $\zeta_n \mapsto \zeta_n^p$. This recovers the isomorphism $\Gal(\mathbb{Q}(\zeta_n)/\mathbb{Q}) \cong (\mathbb{Z}/n\mathbb{Z})^\times$.

Class field theory thus provides a classification of abelian extensions of \emph{any} number field, thereby solving Hilbert's twelfth problem.

\subsection{Non-Abelian Reciprocity and the Langlands Program}

Class field theory handles only \emph{abelian} extensions. The next question is: what about non-abelian extensions? This is the subject of the \textbf{Langlands program}, proposed by Robert Langlands in 1967.

The Langlands program envisions a vast generalization of class field theory in which Galois representations (not just abelian ones) correspond to \emph{automorphic representations}---certain analytic objects arising from harmonic analysis on adelic groups. When the Galois representation is one-dimensional (i.e., abelian), the Langlands program reduces to class field theory.

Why is non-abelian reciprocity so much harder? The abelian case is essentially one-dimensional: abelian extensions correspond to characters of ideal groups, which are relatively simple objects. For non-abelian extensions, the Galois group is more complicated (e.g., $S_3$, $A_5$, or $\GL_2(\mathbb{F}_\ell)$), and the corresponding automorphic objects are far more intricate. There is no simple analogue of roots of unity or the $j$-function.

A significant achievement came in 1995, when Andrew Wiles proved Fermat's Last Theorem by establishing a special case of the Langlands correspondence: the \textbf{modularity theorem}. This theorem states that every elliptic curve over $\mathbb{Q}$ is \textbf{modular}---it corresponds to a modular form, which is an automorphic form for $\GL_2$. In Galois-theoretic terms, the two-dimensional Galois representations arising from the torsion points of an elliptic curve are automorphic.

The modularity theorem can be viewed as a \textbf{non-abelian reciprocity law}: just as the Artin map tells us how primes split in abelian extensions, the modularity theorem tells us that the $L$-function of an elliptic curve equals the $L$-function of a modular form, giving deep constraints on the arithmetic of the curve. This is a non-abelian analogue of the Kronecker--Weber theorem: elliptic curves over $\mathbb{Q}$ are governed by modular forms, much like abelian extensions of $\mathbb{Q}$ are governed by roots of unity.

The full Langlands program extends this vision to all dimensions and all number fields, but it remains one of the major open problems in mathematics.

\section{Current Status}

\begin{solvedbox}
Hilbert's twelfth problem is \textbf{solved}:
\begin{itemize}
    \item The Kronecker--Weber theorem (completed by Weber, 1886) resolves the case $K = \mathbb{Q}$.
    \item Complex multiplication (Kronecker, Fueter, Hasse, 1900s--) resolves the case of imaginary quadratic fields.
    \item Class field theory (Takagi, Artin, 1920s) provides a complete description of abelian extensions for all number fields, though it does not always realize them via special values of analytic functions in the way Kronecker envisioned.
\end{itemize}
The problem is considered solved, but Kronecker's original vision---generating abelian extensions explicitly by special values of analytic functions for \emph{all} number fields---remains an active area of research in the context of Hilbert's 12th problem's modern formulation.
\end{solvedbox}

\section*{Common Questions}

\begin{description}[font=\bfseries, leftmargin=2em, style=nextline]

\item[Q: What exactly was Kronecker's ``dream''?]\hfill\\
Kronecker's \emph{Jugendtraum} (youthful dream) was that every abelian extension of a number field could be generated by special values of certain analytic functions, in the same way that the abelian extensions of $\mathbb{Q}$ are generated by roots of unity $e^{2\pi i/n}$. For $\mathbb{Q}$, the function is the exponential function. For imaginary quadratic fields, the functions are elliptic modular functions. For general number fields, Kronecker believed analogous automorphic functions would do the job. Class field theory ultimately solved the classification problem algebraically, but the explicit analytic construction Kronecker envisioned is fully realized only for $\mathbb{Q}$ and imaginary quadratic fields.

\item[Q: Why cyclotomic fields? What's so special about roots of unity?]\hfill\\
Roots of unity are notable because they combine three properties: (1) they are the values of a simple analytic function $e^{2\pi i z}$ at rational arguments; (2) they satisfy algebraic relations (the cyclotomic polynomials) that are highly structured; and (3) their Galois groups are exactly $(\mathbb{Z}/n\mathbb{Z})^\times$, which are abelian groups. The Kronecker--Weber theorem shows that these three properties are enough to generate \emph{all} abelian extensions of $\mathbb{Q}$. In other words, roots of unity are the ``atomic building blocks'' for the abelian part of the absolute Galois group of $\mathbb{Q}$.

\item[Q: Is class field theory still an active area of research?]\hfill\\
The basic theorems of class field theory (Takagi, Artin, 1920s) are well established and are now part of the standard toolkit in number theory. However, research continues in several directions: \emph{explicit class field theory} (constructing abelian extensions with specific analytic functions for fields beyond $\mathbb{Q}$ and imaginary quadratics), \emph{anabelian geometry} (extending class field theory to non-abelian settings), and the \emph{Langlands program} (the vast generalization of class field theory to all Galois representations). Class field theory is also applied actively in arithmetic geometry, cryptography, and the study of elliptic curves.

\item[Q: What is the Kronecker--Weber theorem in simple terms?]\hfill\\
The Kronecker--Weber theorem says: every finite abelian extension of $\mathbb{Q}$ is contained in some cyclotomic field $\mathbb{Q}(\zeta_n)$. In plain language, if you take any field $L$ that is a finite extension of $\mathbb{Q}$ with an abelian Galois group, then all the numbers in $L$ can be expressed using radicals, specifically roots of unity. For example, $\mathbb{Q}(\sqrt{2})$ is abelian and sits inside $\mathbb{Q}(\zeta_8)$, and $\mathbb{Q}(\sqrt{5})$ sits inside $\mathbb{Q}(\zeta_5)$. The theorem tells us that the maximal abelian extension of $\mathbb{Q}$ is simply the union of all cyclotomic fields.

\item[Q: What are elliptic functions and why do they appear in Kronecker's dream?]\hfill\\
Elliptic functions are doubly periodic meromorphic functions on $\mathbb{C}$---they satisfy $f(z+\omega_1) = f(z+\omega_2) = f(z)$ for two independent periods $\omega_1, \omega_2$. They are the natural generalization of the exponential function $e^{2\pi i z}$ (which is singly periodic) to two dimensions. For imaginary quadratic fields $K$, the period lattice can be chosen so that its endomorphism ring is exactly the ring of integers of $K$ (this is the ``complex multiplication'' property). The values of elliptic functions at division points of this lattice then generate abelian extensions of $K$, exactly as $e^{2\pi i/n}$ generates cyclotomic extensions of $\mathbb{Q}$.

\item[Q: What does ``complex multiplication'' mean?]\hfill\\
An elliptic curve over $\mathbb{C}$ can be written as $\mathbb{C}/\Lambda$ for a lattice $\Lambda$. Normally, the only endomorphisms (structure-preserving maps from the curve to itself) are multiplication by integers. When the lattice $\Lambda$ is such that there are \emph{extra} endomorphisms---multiplication by some algebraic integer not in $\mathbb{Z}$---the curve is said to have \emph{complex multiplication}. This happens precisely when $\Lambda$ is a lattice in an imaginary quadratic field. The extra endomorphisms give the curve a richer arithmetic structure, which is why the resulting elliptic functions generate abelian extensions of the imaginary quadratic field in Kronecker's theory.

\item[Q: How does the Langlands program relate to Kronecker's Jugendtraum?]\hfill\\
The Langlands program can be seen as the \emph{non-abelian} generalization of Kronecker's dream. Kronecker wanted to generate abelian extensions using special functions. The Langlands program seeks a correspondence between $n$-dimensional Galois representations (which need not be abelian) and automorphic forms on $\mathrm{GL}_n$. When $n=1$, this reduces to class field theory and hence to Kronecker's abelian setting. When $n > 1$, the correspondence gives a way to understand non-abelian extensions using analytic objects (automorphic forms)---a direct analogue of Kronecker's vision, but vastly more complex. Wiles's proof of Fermat's Last Theorem (a $\mathrm{GL}_2$ case) is arguably the greatest success of this program so far.

\end{description}

\section*{Exercises}

\begin{enumerate}

    \item \textbf{Galois groups of cyclotomic fields.}
    \begin{enumerate}
        \item Prove that $\Gal(\mathbb{Q}(\zeta_n)/\mathbb{Q}) \cong (\mathbb{Z}/n\mathbb{Z})^\times$.
        \item Compute the Galois groups of $\mathbb{Q}(\zeta_5)$, $\mathbb{Q}(\zeta_7)$, and $\mathbb{Q}(\zeta_{12})$.
        \item Determine which of these groups are cyclic.
    \end{enumerate}

    \item \textbf{The Kronecker--Weber theorem in practice.}
    \begin{enumerate}
        \item Show that $\sqrt{2} \in \mathbb{Q}(\zeta_8)$ by expressing $\sqrt{2}$ as $\zeta_8 + \zeta_8^{-1}$.
        \item Show that $\sqrt{3} \in \mathbb{Q}(\zeta_{12})$.
        \item Find an $n$ such that $\mathbb{Q}(\sqrt{-3}) \subseteq \mathbb{Q}(\zeta_n)$.
        \item More generally, prove that for any squarefree $d$, $\mathbb{Q}(\sqrt{d}) \subseteq \mathbb{Q}(\zeta_{|d|})$ when $d$ is odd and $d \equiv 1 \pmod{4}$; otherwise $\mathbb{Q}(\sqrt{d}) \subseteq \mathbb{Q}(\zeta_{4|d|})$.
    \end{enumerate}

    \item \textbf{Finding cyclotomic fields for quadratic fields.}
    \begin{enumerate}
        \item Find the smallest $n$ such that $\mathbb{Q}(\sqrt{-1}) \subseteq \mathbb{Q}(\zeta_n)$. (Answer: $n=4$.)
        \item Find the smallest $n$ such that $\mathbb{Q}(\sqrt{-2}) \subseteq \mathbb{Q}(\zeta_n)$. (Hint: consider $\zeta_8$.)
        \item Find the smallest $n$ such that $\mathbb{Q}(\sqrt{6}) \subseteq \mathbb{Q}(\zeta_n)$.
        \item Show that $\mathbb{Q}(\sqrt{-5}) \subseteq \mathbb{Q}(\zeta_{20})$ but not $\mathbb{Q}(\zeta_5)$.
    \end{enumerate}

    \item \textbf{Abelian vs.\ non-abelian extensions.}
    \begin{enumerate}
        \item The splitting field of $x^3 - 2$ over $\mathbb{Q}$ is $\mathbb{Q}(\sqrt[3]{2}, \zeta_3)$. Compute its Galois group and show it is not abelian.
        \item Why does the Kronecker--Weber theorem not apply to this extension?
        \item Let $L$ be the splitting field of $x^5 - 1$. Show $L = \mathbb{Q}(\zeta_5)$ and compute its Galois group.
    \end{enumerate}

    \item \textbf{Subfields of cyclotomic fields.}
    \begin{enumerate}
        \item The field $\mathbb{Q}(\zeta_5)$ has degree $4$ over $\mathbb{Q}$. Find all its subfields by analyzing the subgroup structure of $(\mathbb{Z}/5\mathbb{Z})^\times \cong C_4$.
        \item Repeat for $\mathbb{Q}(\zeta_7)$, which has degree $6$.
        \item Show that $\mathbb{Q}(\sqrt{5})$ is the unique quadratic subfield of $\mathbb{Q}(\zeta_5)$.
    \end{enumerate}

    \item \textbf{The Artin map.}
    \begin{enumerate}
        \item For the cyclotomic extension $\mathbb{Q}(\zeta_n)/\mathbb{Q}$, describe the Artin map explicitly. (Hint: the Artin map sends a prime $p \nmid n$ to the automorphism $\zeta_n \mapsto \zeta_n^p$.)
        \item Verify that the kernel of the Artin map consists of principal ideals generated by integers congruent to $1$ modulo $n$.
        \item How does this relate to the isomorphism $\Gal(\mathbb{Q}(\zeta_n)/\mathbb{Q}) \cong (\mathbb{Z}/n\mathbb{Z})^\times$?
    \end{enumerate}

    \item \textbf{The Artin map for quadratic fields.}
    \begin{enumerate}
        \item Let $d$ be a squarefree integer. Show that for a prime $p$ not dividing $d$, the Artin map sends $p$ to the automorphism $\sqrt{d} \mapsto \left(\frac{d}{p}\right)\sqrt{d}$, where $\left(\frac{d}{p}\right)$ is the Legendre symbol.
        \item Use quadratic reciprocity to show that the Artin map factors through the ray class group modulo $|d|$ (or $4|d|$).
        \item Verify this explicitly for $d = 5$ and $p = 11, 13, 17, 19$.
    \end{enumerate}

    \item $\star$ \textbf{Conductors of quadratic fields.}
    The \textbf{conductor} of an abelian extension $L/\mathbb{Q}$ is the smallest $n$ such that $L \subseteq \mathbb{Q}(\zeta_n)$.
    \begin{enumerate}
        \item Find the conductors of $\mathbb{Q}(\sqrt{2})$, $\mathbb{Q}(\sqrt{3})$, $\mathbb{Q}(\sqrt{5})$, and $\mathbb{Q}(\sqrt{-1})$.
        \item Conjecture a formula for the conductor of $\mathbb{Q}(\sqrt{d})$ in terms of $d$.
        \item (Harder) Show that the conductor of $\mathbb{Q}(\sqrt{d})$ for squarefree $d$ is $|d|$ if $d \equiv 1 \pmod{4}$, and $4|d|$ otherwise.
    \end{enumerate}

    \item $\star$ \textbf{Complex multiplication.}
    \begin{enumerate}
        \item The $j$-invariant evaluated at $\tau = i$ is $j(i) = 1728$. Verify that $j(i)$ is an integer, and show that $\mathbb{Q}(j(i)) = \mathbb{Q}$.
        \item For $\tau = e^{2\pi i/3}$, we have $j(e^{2\pi i/3}) = 0$. What does this tell us about the abelian extensions of $\mathbb{Q}(\sqrt{-3})$?
        \item Explain why the values $j(\tau)$ for $\tau$ in an imaginary quadratic field $K$ are algebraic integers, and why adjoining them to $K$ yields abelian extensions.
    \end{enumerate}

    \item $\star\star$ \textbf{Non-abelian reciprocity and elliptic curves.}
    \begin{enumerate}
        \item Look up the statement of the modularity theorem. Why is it called a ``non-abelian reciprocity law''?
        \item The elliptic curve $E: y^2 + y = x^3 - x$ has conductor $37$. Find the corresponding modular form and compute its first few Fourier coefficients.
        \item Explain how the Langlands program generalizes the Artin reciprocity law.
    \end{enumerate}

\end{enumerate}
